%%
%% winter-lecture26.tex
%% 
%% Made by Alex Nelson <pqnelson@gmail.com>
%% Login   <alex@lisp>
%% 
%% Started on  2026-03-05T09:48:03-0800
%% Last update 2026-03-05T09:48:03-0800
%% 

\lecture{}

We used the following fact last time (but we did not prove it):

\begin{proposition}
If $\{\varphi_{k}(x)\}$ is an orthonormal basis for $L^{2}(\RR^{d})$,
then $\Phi_{k,\ell}(x,y)=\varphi_{k}(x)\varphi_{\ell}(y)$ is an
orthonormal basis for $L^{2}(\RR^{d}\times\RR^{d})$.
\end{proposition}

\begin{proof}
We need to prove orthonormality and completeness.

\textsc{Claim 1: Orthonormality}. We can find this directly by
\begin{subequations}
  \begin{align}
\iint\Phi_{k,\ell}(x,y)\overline{\Phi_{m,n}(x,y)}\,\D x\D y
&=\iint\varphi_{k}(x)\varphi_{\ell}(y)\overline{\varphi_{m}(x)\varphi_{n}(y)}\,\D x\D y\\
&=\int\varphi_{k}(x)\overline{\varphi_{m}(x)}\,\D x\int\varphi_{\ell}(y)\overline{\varphi_{n}(y)}\,\D y\mbox{ by Fubini}\\
&=\delta_{m,k}\delta_{n,\ell}
  \end{align}
\end{subequations}
by orthonormality of the $\varphi_{k}$.

\textsc{Claim 2: Completeness}. Let $f\in L^{2}(\RR^{d}\times\RR^{d})$
be orthogonal to all $\Phi_{k,\ell}$. We want to show $f=0$ almost
everywhere.
For each $\ell$, define $f_{\ell}\in L^{2}(\RR^{d})$ by
\begin{equation}
f_{\ell}(x)=\int_{\RR^{d}}f(x,y)\overline{\varphi_{\ell}(y)}\,\D y.
\end{equation}
By assumption, for any $k$,
\begin{equation}
\langle f_{\ell},\varphi_{k}\rangle_{L^{2}(\RR^{d})}=0
\end{equation}
and since $\{\varphi_{k}\}$ is a basis for $L^{2}(\RR^{d})$ this means
that $f_{\ell}(x)=0$ for almost all $x\in\RR^{d}$. So then $y\mapsto f(x,y)$
is orthogonal to every $\varphi_{\ell}(y)$, and since $\{\varphi_{\ell}(y)\}$
is a basis, this means
\begin{equation}
f(x,y)=0\mbox{ for almost every }(x,y)\in\RR^{d}\times\RR^{d}.
\end{equation}
Hence $\Phi_{k,\ell}$ is complete.
\end{proof}

\begin{theorem}[Uniform boundedness principle]
Let $X$ be a Banach space, let $Y$ be a normed vector space, let
$\{T_{\alpha}\colon X\to Y\}_{\alpha\in A}$ be a family of continuous
linear operators.

If we have pointwise boundedness ($\forall x\in X\ldotp\sup_{\alpha\in A}\|T_{\alpha}x\|_{Y}<\infty$),
then we have \define{Uniform Boundedness}, i.e., $\sup_{\alpha\in A}\|T_{\alpha}\|_{\text{op}}<\infty$.
\end{theorem}

\begin{proof}
Let
\begin{equation}
E_{n}:=\{x\in X\mid\forall\alpha\in A\ldotp\|T_{\alpha}x\|\leq n\}.
\end{equation}
Then for any $x\in X$, there is a real number $M_{x}\in\RR$ such that
\begin{equation}
\sup_{\alpha\in A}\|T_{\alpha}x\|=M_{x}
\end{equation}
by hypothesis. If $n\geq M_{x}$, then $x\in E_{n}$. In other words,
every $x\in X$ eventually lands in $E_{n}$. This gives us a decomposition
\begin{equation}
X=\bigcup_{n=1}^{\infty}E_{n}.
\end{equation}
Great.

Now we claim these $E_{n}$ are closed subsets. We see
\begin{subequations}
  \begin{align}
E_{n} &= \bigcap_{\alpha\in A}\{x\in X\mid \|T_{\alpha}x\|\leq n\}\\
&=\bigcap_{\alpha\in A}T_{\alpha}^{-1}(\closure{B_{Y}(0,r=n)})
  \end{align}
\end{subequations}
where $B_{Y}(y_{0},r)$ is the open ball in $Y$ centered at $y_{0}$ of
radius $r>0$; now, since $T_{\alpha}$ is continuous, each
$T_{\alpha}^{-1}(\closure{B_{Y}(0,r=n)})$ is closed (the preimage
under a continuous map of a closed set is closed), and the arbitrary
intersection of closed sets is closed. Therefore $E_{n}$ is closed.

Now, using the Baire category theorem, $X$ cannot be meager (i.e., $X$
cannot be written as a countable union of nowhere dense sets). Then it
is impossible for all the $E_{n}$ to be nowhere dense. We saw
$E_{n}=\closure{E_{n}}$ are closed subsets. For a closed set to be
nowhere dense, it must have a empty interior. Since $E_{n}\subset E_{n+1}$
is ascending, there must exist a positive integer $N\in\NN$ such that
\begin{equation}
\Interior(E_{N})\neq\emptyset
\end{equation}
since not all $E_{n}$ are nowhere dense.

This means there exists a solid open ball
\begin{equation}
B_{X}(x_{0},r)\propersubset E_{N}
\end{equation}
for some $x_{0}\in X$ and $r>0$. (We want to move this ball to the origin.)

If $x\in B_{X}(0,r)$, then
\begin{equation}
x+x_{0}\in B_{X}(x_{0},r)\propersubset E_{N},
\end{equation}
and by definition of $E_{N}$, we have bounds
\begin{subequations}
  \begin{align}
\|T_{\alpha}(x+x_{0})\|\leq N\\
\intertext{and}
\|T_{\alpha}x_{0}\|\leq N.
  \end{align}
\end{subequations}
Then by linearity of $T_{\alpha}$ and the triangle inequality, we have
\begin{subequations}
  \begin{align}
\|T_{\alpha}x\|_{Y}
&=\|T_{\alpha}(x+x_{0}-x_{0})\|_{Y}\\
&=\|T_{\alpha}(x_{0}+x+(-x_{0}))\|_{Y}\\
&\leq\|T_{\alpha}(x_{0}+x)\|_{Y}+\|T_{\alpha}(-x_{0}))\|_{Y}\mbox{ by triangle ineq.}\\
&\leq N + N=2N.
  \end{align}
\end{subequations}
But since $x\in B_{X}(0,r)$ we let $u$ be any unit vector $\|u\|\leq1$
such that $ru\in B_{X}(0,r)$, then we must have (by the previous equation),
\begin{equation}
\|T_{\alpha}(ru)\|_{Y}\leq2N
\end{equation}
and by linearity
\begin{equation}
\|T_{\alpha}(ru)\|_{Y}=|r|\|T_{\alpha}(u)\|_{Y}\leq 2N
\end{equation}
so divide both sides by $r=|r|$ (since $r>0$),
\begin{equation}
\|T_{\alpha}(u)\|_{Y}\leq\frac{2N}{r}.
\end{equation}
But since this is for any $u\in B_{X}(0,1)$ and fixed $N\in\NN$ and $r>0$,
and we see the left-hand side is just the operator norm
\begin{equation}
\|T_{\alpha}\|_{\text{op}}\leq\frac{2N}{r}
\end{equation}
and the right-hand side is ``some [positive finite] constant''. This
means the operator norm of $T_{\alpha}$ is bounded for all $\alpha\in A$.
Hence the claim.
\end{proof}

\begin{observation}
In the proof, completeness of $X$ gave us the ball, linearity of
$T_{\alpha}$ let us move it.
\end{observation}

\begin{definition}
A map $f\colon X\to Y$ (between topological spaces) is \define{Open} if for every open subset
$U\subset X$ its image $f(U)\subset Y$ is open.
\end{definition}

\begin{recall}
We say $f\colon X\to Y$ is continuous if for every open subset
$V\subset Y$ its preimage $f^{-1}(V)\subset X$ is open.
\end{recall}

\begin{example}[Continuous but not open]
Consider $f\colon\RR\to\RR$ given by $f(x)=x^{2}$. Then $f$ maps the
open interval $(-1,1)$ to the clopen interval $[0,1]$. But $[0,1)$ is
not an open subset of $\RR$, so $f$ cannot possibly be open.
\end{example}

\begin{example}[Open but not continuous]
Consider the function $g\colon\RR^{2}\to\RR$ given by
$g(x,y)=\sgn(y)+x$ where
\begin{equation}
\sgn(y)=\begin{cases}-1 & \mbox{if }y<0\\
0 & \mbox{if }y=0,\\
+1 & \mbox{if }y>0.
\end{cases}
\end{equation}
Then $g$ is obviously not continuous in a neighborhood where $y=0$,
because there is a discontinuous jump from $\sgn(y)$.

But we claim $g$ is an open map. Let $U\subset\RR^{2}$ be any open
set, let $z\in g(U)$. Then
\begin{equation}
z=x_{0}+\sgn(y_{0})
\end{equation}
for some $(x_{0},y_{0})\in U$. Then there is some $\varepsilon>0$ such
that we have some line in $U$ containing $(x_{0},y_{0})$,
\begin{equation}
L:=(x_{0}-\varepsilon,x_{0}+\varepsilon)\times\{y_{0}\}\subset U.
\end{equation}
On $L$, we see $\sgn(y)=\sgn(y_{0})$ is constant, so
\begin{equation}
g|_{L}(x,y)=x+\sgn(y_{0}),
\end{equation}
and
\begin{equation}
\begin{split}
g(L)&=(x_{0}+\sgn(y_{0})-\varepsilon,x_{0}+\sgn(y_{0})+\varepsilon)\\
&=(z-\varepsilon,z+\varepsilon)
\end{split}
\end{equation}
This means $g(L)$ is an open subset of $\RR$, specifically an open
neighborhood of $z\in g(U)$ and $g(L)\subset g(U)$. We let $z\in g(U)$
be arbitrary and we can always find such lines $L$ for which
$g(L)\subset g(U)$ is open, which means $g(U)$ must be open.
\end{example}

\begin{theorem}[Open Mapping Theorem]
Let $X$ and $Y$ both be Banach spaces.
Let $T\colon X\to Y$ be a bounded linear operator.
If $T$ is surjective, then $T$ is an open map.
\end{theorem}

\begin{remark}
Suffices to show $T(B_{X}(0,r=1))\supset B_{Y}(0,r)$ for some
$r>0$. Then lienarity allows us to scale $B_{X}(0,1)$ and translate
it, boundedness completes the proof.
\end{remark}

\begin{lemma}[Core Lemma]If $T\colon X\to Y$ is a surjective bounded linear operator
between Banach spaces $X$ and $Y$, then there exists some $r>0$ such
that
\begin{equation}
\closure{T(B_{x}(0,1))}\supset B_{Y}(0,r).
\end{equation}
\end{lemma}

\begin{proof}
\textsc{Step 1:} Since $T$ is surjective, every $y\in Y$ is hit by
some $x\in X$. Since $X$ has a finite norm, $x$ lives in some open
ball $B_{X}(0,n)$ for some $n\in\NN$. So
\begin{equation}
Y=\bigcup^{\infty}_{n=1}T(B_{X}(0,n))
\end{equation}
[$Y\subset\bigcup T(B(0,n))$ by surjectivity and $Y\supset\bigcup T(B(0,n))$
  because duh]. But
\begin{equation}
Y=\bigcup_{n\in\NN}T(B_{n})\subset\bigcup_{n\in\NN}\closure{T(B_{n})}\subset Y
\end{equation}
where $B_{n}=B_{x}(0,n)$.

\textsc{Step 2}: By hypothesis, $Y$ is Banach, so it is a
\emph{complete} metric space, which is a countable union of closed
sets. By  the Baire category theorem, at least one of these sets has a
nonempty interior: there exits some $N\in\NN$ such that
$\closure{T}(B_{N})$ contains some open ball $B_{Y}(y_{0},r)\subset\closure{T}(B_{N})$.

\textsc{Step 3: Scaling}. By linearity,
\begin{equation}
\closure{T(B_{N})}=N\closure{T(B_{1})},
\end{equation}
then we can divide by $N$ to find
\begin{equation}
\closure{T(B_{1})}\supset B_{Y}(y_{1},r')
\end{equation}
where $y_{1}=y_{0}/N$ and $r'=r/N$.

\textsc{Step 4:} We claim $B_{Y}(0,r'/2)\subset\closure{T(B_{1})}$.
\end{proof}